Question~9.1 of \emph{Directed Recovery Structures in Tape Computations} asks whether every storage code on a recovery
relation of bounded incidence degree with the paired-distance property has rate at most $1/2+o(1)$; the fifth paper of
the series proposed it as a machine-free route to lower bounds for linear-time matrix transposition on multitape
machines. We show that the answer is negative. For the prime $p=127$ we construct recovery relations of incidence degree at
most an explicit, astronomically large bound $\Delta<127^{8\,791\,448}$, on which the paired distance of every coordinate
tends to infinity, carrying $\F_{127^2}$-linear storage codes of rate at least $8100/16130>0.5021$; writing each symbol in
$14$ bits gives binary storage codes of rate above $0.5013$. For large primes inert in $\Q(\sqrt5)$ the guaranteed rates
approach $4-2\sqrt3\approx0.536$. Forgetting the directions of the relations gives graphs of bounded degree in which every odd
cycle is long, carrying the same codes, so the fourth paper's undirected crux (its Question~6.1, on locally bipartite graphs of bounded
degree) has a negative answer as well. The relations are twisted quotients of the Bruhat--Tits tree of $\PGL_2(\Q_{p^2})$, and the codes
are Tanner codes with Reed--Solomon local codes. Their large dimension comes from Riemann--Roch on Hilbert modular surfaces
in characteristic $p$, through the Goren--Oort divisor, whose dual graph is the quotient of the tree. The combinatorial
ingredients are lift selections of degeneracy $(2-\sqrt3)\delta+O(\sqrt{\delta\log n})$ in double covers of $\delta$-regular
graphs on $n$ vertices, and private checks built by peeling, which make the degree bound explicit.

The combinatorics is formalized in Lean~4 with Mathlib, on the standard axioms only: the deterministic part of the
lift-selection theorem, the peeling lemma, the descent to quotients and its symmetrization, the binary blow-ups, the
numerical facts used in the proofs, and the deduction of the negation of the eleventh paper's formal statement of
Question~9.1 from a tower of quotients. Not formalized are the probabilistic step of the lift-selection theorem, the
covering-theoretic passage to the tree, the count of balls in the tree, the structure of the icosian group and its
congruence subgroups with the girth bound, the arithmetic dimension bound, which uses Riemann--Roch, partial Hasse
invariants and the uniformization of the supersingular locus, and the extension to all large primes.
